\begin{abstract}
When two LLM agents want different things from a shared decision, how should each handle the conflict? Organizational psychology names five conflict-handling strategies (competing, collaborating, compromising, accommodating, avoiding) and claims that each fits some situations. We test these claims in a parametric multi-issue bargaining environment with an exact efficiency frontier, giving each strategy as behavioural rules whose compliance is checked from the moves, over about 20{,}000 games: the full strategy matrix, ten situational conditions, three framings, instances derived from real code conflicts, two published benchmarks, four models and cross-model pairs. In the matrix, collaborating is the best response to almost every partner, competing pays only by transfer and only against a partner that cannot push back, and compromising and avoiding destroy value. The situation moves the map, but not as the textbook prescribes, and a price on tokens moves the answer to the cheap strategies early. Across 29 settings what transfers is what not to do: competing, compromising and avoiding pairs fall below the default pair in 80 of 81 comparisons, and the soft strategies lose when adopted alone. What does not transfer is who keeps collaborating's gain: the rule raises the joint payoff with almost every partner, and its user keeps the gain only against a partner that cannot find the frontier alone, as pairs of different models show directly. The environment and all transcripts are released.
\end{abstract}

